% F085｜八类模型从哪里来、能解决什么问题。
\FReserve{7.5cm}
\begin{explanationbox}
\textbf{\textcolor{CExplanation}{1 指数型：为什么“加法变乘法”？}}
% F085｜真实数学图；代表函数不构成抽象解的唯一性证明。
\begin{center}
\begin{tikzpicture}
\begin{axis}[width=12.4cm,height=5.1cm,axis lines=middle,
 xmin=-1.4,xmax=3.5,ymin=-0.5,ymax=9.5,
 xtick={-1,0,1,2,3},xticklabels=\empty,ytick={0,2,4,8},xlabel={$t$},ylabel={$f(t)$},
 title={指数曲线 $f(t)=2^t$},title style={font=\small,yshift=3pt},
 tick label style={font=\small},label style={font=\small},
 ylabel style={at={(axis description cs:0,1)},anchor=south east,xshift=-18pt},
 clip=false,grid=major,grid style={black!8},enlargelimits=false,
 every axis plot/.append style={thick},samples=161]
\addplot[CExplanation,very thick,domain=-1.2:3.15] {2^x};
\draw[dashed,black!35] (axis cs:1,0) |- (axis cs:0,2);
\draw[dashed,black!35] (axis cs:2,0) |- (axis cs:0,4);
\draw[dashed,black!35] (axis cs:3,0) |- (axis cs:0,8);
\addplot[only marks,mark=*,mark size=2pt,orange!80!black] coordinates {(1,2) (2,4) (3,8)};

\node[font=\small,anchor=north,yshift=-3pt] at (axis cs:-1,-0.5) {$-1$};
\node[font=\small,anchor=north,yshift=-3pt] at (axis cs:1,-0.5) {$1$};
\node[font=\small,anchor=north,yshift=-3pt] at (axis cs:2,-0.5) {$2$};
\node[font=\small,anchor=north,yshift=-3pt] at (axis cs:3,-0.5) {$3$};

\end{axis}
\end{tikzpicture}
\par\small 横坐标 $1+2=3$，对应高度从 $2,4$ 相乘得到 $8$；曲线经过 $(0,1)$。
\end{center}


图上 $f(1)=2,f(2)=4,f(3)=8$：横坐标相加 $1+2=3$，高度相乘 $2\times4=8$。一般地，$2^{x+y}=2^x2^y$。
\[
\boxed{f(x+y)=f(x)f(y)}.
\]
\textbf{考试中用它：}看到 $f(2x)$、$f(3x)$ 或 $f(m+n)$，直接拆输入；已知 $f(1)$ 就能求很多整数点的值。

\textbf{边界：}该式也允许 $f\equiv0$；不能看到等式就把 $f(0)=1$ 当作无条件结论。在全体实数上连续且不恒为零时，才可进一步确定为 $a^x$（$a>0$）。
\end{explanationbox}

\FReserve{7.5cm}
\begin{explanationbox}
\textbf{\textcolor{CExplanation}{2 对数型：为什么“乘法变加法”？}}
% F085｜真实数学图；代表函数不构成抽象解的唯一性证明。
\begin{center}
\begin{tikzpicture}
\begin{axis}[width=12.4cm,height=5.1cm,axis lines=middle,
 xmin=-0.3,xmax=8.9,ymin=-2.7,ymax=3.8,
 xtick={1,2,4,8},xticklabels=\empty,ytick={-2,-1,0,1,2,3},xlabel={$t$},ylabel={$f(t)$},
 title={对数曲线 $f(t)=\log_2t$},title style={font=\small,yshift=3pt},
 tick label style={font=\small},label style={font=\small},
 ylabel style={at={(axis description cs:0,1)},anchor=south east,xshift=-18pt},
 clip=false,grid=major,grid style={black!8},enlargelimits=false,
 every axis plot/.append style={thick},samples=161]
\addplot[CExplanation,very thick,domain=0.2:8.3] {ln(x)/ln(2)};
\draw[dashed,black!35] (axis cs:2,0) |- (axis cs:0,1);
\draw[dashed,black!35] (axis cs:4,0) |- (axis cs:0,2);
\draw[dashed,black!35] (axis cs:8,0) |- (axis cs:0,3);
\addplot[only marks,mark=*,mark size=2pt,orange!80!black] coordinates {(2,1) (4,2) (8,3)};

\node[font=\small,anchor=north,yshift=-3pt] at (axis cs:1,-2.7) {$1$};
\node[font=\small,anchor=north,yshift=-3pt] at (axis cs:2,-2.7) {$2$};
\node[font=\small,anchor=north,yshift=-3pt] at (axis cs:4,-2.7) {$4$};
\node[font=\small,anchor=north,yshift=-3pt] at (axis cs:8,-2.7) {$8$};

\end{axis}
\end{tikzpicture}
\par\small 输入从 $2$ 到 $4$ 再到 $8$，每次乘 $2$；高度从 $1$ 到 $2$ 再到 $3$，每次加 $1$。
\end{center}


图上横坐标相乘 $2\cdot4=8$，高度相加 $1+2=3$，即
\[
\log_2(2\cdot4)=\log_2 2+\log_2 4=1+2.
\]
因而熟悉的对数函数满足
\[\boxed{f(xy)=f(x)+f(y)}\qquad(x,y>0).\]
\textbf{考试中用它：}把 $8=2^3$ 看成三个 $2$ 相乘，便有 $f(8)=3f(2)$。

\textbf{注意：}这里的输入是正数；$f(1)=0$，因为 $f(1)=f(1\cdot1)=2f(1)$。
\end{explanationbox}

\FReserve{7.5cm}
\begin{explanationbox}
\textbf{\textcolor{CExplanation}{3 幂函数型：“乘法还是乘法”}}
% F085｜真实数学图；代表函数不构成抽象解的唯一性证明。
\begin{center}
\begin{tikzpicture}
\begin{axis}[width=12.4cm,height=5.1cm,axis lines=middle,
 xmin=-0.3,xmax=6.7,ymin=-2,ymax=41,
 xtick={0,1,2,3,6},xticklabels=\empty,ytick={0,4,9,36},xlabel={$t$},ylabel={$f(t)$},
 title={幂函数曲线 $f(t)=t^2$（这里取 $t>0$）},title style={font=\small,yshift=3pt},
 tick label style={font=\small},label style={font=\small},
 ylabel style={at={(axis description cs:0,1)},anchor=south east,xshift=-18pt},
 clip=false,grid=major,grid style={black!8},enlargelimits=false,
 every axis plot/.append style={thick},samples=161]
\addplot[CExplanation,very thick,domain=0:6.2] {x^2};
\draw[dashed,black!35] (axis cs:2,0) |- (axis cs:0,4);
\draw[dashed,black!35] (axis cs:3,0) |- (axis cs:0,9);
\draw[dashed,black!35] (axis cs:6,0) |- (axis cs:0,36);
\addplot[only marks,mark=*,mark size=2pt,orange!80!black] coordinates {(2,4) (3,9) (6,36)};

\node[font=\small,anchor=north,yshift=-3pt] at (axis cs:1,-2) {$1$};
\node[font=\small,anchor=north,yshift=-3pt] at (axis cs:2,-2) {$2$};
\node[font=\small,anchor=north,yshift=-3pt] at (axis cs:3,-2) {$3$};
\node[font=\small,anchor=north,yshift=-3pt] at (axis cs:6,-2) {$6$};

\end{axis}
\end{tikzpicture}
\par\small 横坐标相乘 $2\cdot3=6$，高度也相乘 $4\cdot9=36$；$(0,0)$ 仅是边界参照。
\end{center}


取 $f(x)=x^2$，则
\[
 f(2\cdot3)=6^2=36=4\times9=f(2)f(3).
\]
因此 $\boxed{f(xy)=f(x)f(y)}$ 是幂型模型的显著特征。

\textbf{与对数的对比：}在相同的正数输入上，若 $f(2)=3$，则
\[
\begin{array}{ll}
\text{对数型：}&f(8)=3f(2)=9;\\
\text{幂函数型：}&f(8)=[f(2)]^3=27.
\end{array}
\]
\textbf{注意：}不恒为零的解在正数范围内有 $f(1)=1$、$f(x)>0$；常见的 $x^a$ 是连续情况下的代表。
\end{explanationbox}

\FReserve{7.5cm}
\begin{explanationbox}
\textbf{\textcolor{CExplanation}{4 正比例型：“加法还是加法”}}
% F085｜真实数学图；代表函数不构成抽象解的唯一性证明。
\begin{center}
\begin{tikzpicture}
\begin{axis}[width=12.4cm,height=5.1cm,axis lines=middle,
 xmin=-2.3,xmax=3.5,ymin=-7,ymax=10.5,
 xtick={-2,-1,0,1,2,3},xticklabels=\empty,ytick={-6,-3,0,3,6,9},xlabel={$t$},ylabel={$f(t)$},
 title={正比例直线 $f(t)=3t$},title style={font=\small,yshift=3pt},
 tick label style={font=\small},label style={font=\small},
 ylabel style={at={(axis description cs:0,1)},anchor=south east,xshift=-18pt},
 clip=false,grid=major,grid style={black!8},enlargelimits=false,
 every axis plot/.append style={thick},samples=161]
\addplot[CExplanation,very thick,domain=-2:3.1] {3*x};
\draw[dashed,black!35] (axis cs:1,0) |- (axis cs:0,3);
\draw[dashed,black!35] (axis cs:2,0) |- (axis cs:0,6);
\draw[dashed,black!35] (axis cs:3,0) |- (axis cs:0,9);
\draw[dashed,black!35] (axis cs:-1,0) |- (axis cs:0,-3);
\addplot[only marks,mark=*,mark size=2pt,orange!80!black] coordinates {(1,3) (2,6) (3,9) (-1,-3)};

\node[font=\small,anchor=north,yshift=-3pt] at (axis cs:-2,-7) {$-2$};
\node[font=\small,anchor=north,yshift=-3pt] at (axis cs:-1,-7) {$-1$};
\node[font=\small,anchor=north,yshift=-3pt] at (axis cs:1,-7) {$1$};
\node[font=\small,anchor=north,yshift=-3pt] at (axis cs:2,-7) {$2$};
\node[font=\small,anchor=north,yshift=-3pt] at (axis cs:3,-7) {$3$};

\end{axis}
\end{tikzpicture}
\par\small 输入相加，高度相加；点 $(1,3)$ 与 $(-1,-3)$ 关于原点对称。
\end{center}


图上直线 $f(x)=3x$ 满足 $f(1+2)=9=3+6=f(1)+f(2)$。
\[\boxed{f(x+y)=f(x)+f(y)}.\]

\textbf{为什么重要？}这一式马上给出 $f(0)=0$、$f(-x)=-f(x)$，以及 $f(nx)=nf(x)$（整数 $n$）。

\textbf{但不要急着说 $f(x)=kx$：}在有连续性（或单调性等）条件时，才能按正比例函数确定解析式。单凭关系式，先做代入计算更稳妥。
\end{explanationbox}

\FReserve{7.5cm}
\begin{explanationbox}
\textbf{\textcolor{CExplanation}{5 余弦型：为什么左右两项相加？}}
% F085｜真实数学图；代表函数不构成抽象解的唯一性证明。
\begin{center}
\begin{tikzpicture}
\begin{axis}[width=12.4cm,height=5.1cm,axis lines=middle,
 xmin=-0.3,xmax=6.8,ymin=-1.5,ymax=1.5,
 xtick={0,1.570796,3.141593,4.712389,6.283185},xticklabels=\empty,ytick={-1,0,1},xlabel={$t$},ylabel={$f(t)$},
 title={余弦曲线 $f(t)=\cos t$},title style={font=\small,yshift=3pt},
 tick label style={font=\small},label style={font=\small},
 ylabel style={at={(axis description cs:0,1)},anchor=south east,xshift=-18pt},
 clip=false,grid=major,grid style={black!8},enlargelimits=false,
 every axis plot/.append style={thick},samples=161]
\addplot[CExplanation,very thick,domain=0:6.28] {cos(deg(x))};
\draw[dashed,black!35] (axis cs:0.5235987755982988,0) |- (axis cs:0,0.8660254037844386);
\draw[dashed,black!35] (axis cs:1.5707963267948966,0) |- (axis cs:0,0);
\addplot[only marks,mark=*,mark size=2pt,orange!80!black] coordinates {(0.523598776,0.866025404) (1.57079633,0)};

\node[font=\small,anchor=north,yshift=-3pt] at (axis cs:1.570796,-1.5) {$\pi/2$};
\node[font=\small,anchor=north,yshift=-3pt] at (axis cs:3.141593,-1.5) {$\pi$};
\node[font=\small,anchor=north,yshift=-3pt] at (axis cs:4.712389,-1.5) {$3\pi/2$};
\node[font=\small,anchor=north,yshift=-3pt] at (axis cs:6.283185,-1.5) {$2\pi$};

\end{axis}
\end{tikzpicture}
\par\small 取 $x\pm y=\pi/2,\pi/6$：两个高度之和 $0+\sqrt3/2=2\cos x\cos y$。横轴单位为弧度。
\end{center}


两角和差公式里，$\sin x\sin y$ 在一个式子中为负、另一个式子中为正；相加后恰好抵消：
\[
\begin{aligned}
\cos(x+y)&=\cos x\cos y-\sin x\sin y,\\
\cos(x-y)&=\cos x\cos y+\sin x\sin y.
\end{aligned}
\]
所以
\[\boxed{\cos(x+y)+\cos(x-y)=2\cos x\cos y}.\]

\textbf{考试中用它：}令 $y=x$，就能把 $f(2x)$ 变成 $f(x)$ 的平方关系（非零解需先推出 $f(0)=1$）。
\end{explanationbox}

\FReserve{7.5cm}
\begin{explanationbox}
\textbf{\textcolor{CExplanation}{6 正切型：加法变成一个分式}}
% F085｜真实数学图；代表函数不构成抽象解的唯一性证明。
\begin{center}
\begin{tikzpicture}
\begin{axis}[width=12.4cm,height=5.1cm,axis lines=middle,
 xmin=-1.4,xmax=2.9,ymin=-3.7,ymax=3.7,
 xtick={-1,0,1,2},xticklabels=\empty,ytick={-3,-1,0,1,3},xlabel={$t$},ylabel={$f(t)$},
 title={正切曲线 $f(t)=\tan t$},title style={font=\small,yshift=3pt},
 tick label style={font=\small},label style={font=\small},
 ylabel style={at={(axis description cs:0,1)},anchor=south east,xshift=-18pt},
 clip=false,grid=major,grid style={black!8},enlargelimits=false,
 every axis plot/.append style={thick},samples=161]
\addplot[CExplanation,very thick,domain=-1.2:1.28] {tan(deg(x))};
\addplot[CExplanation,very thick,domain=1.86:2.7] {tan(deg(x))};
\draw[dashed,orange!80!black] (axis cs:1.5707963267948966,-3.5)--(axis cs:1.5707963267948966,3.5);
\draw[dashed,black!35] (axis cs:0.5235987755982988,0) |- (axis cs:0,0.5773502691896258);
\draw[dashed,black!35] (axis cs:1.0471975511965976,0) |- (axis cs:0,1.7320508075688772);
\addplot[only marks,mark=*,mark size=2pt,orange!80!black] coordinates {(0.523598776,0.577350269) (1.04719755,1.73205081)};

\node[font=\small,anchor=north,yshift=-3pt] at (axis cs:-1,-3.7) {$-1$};
\node[font=\small,anchor=north,yshift=-3pt] at (axis cs:1,-3.7) {$1$};
\node[font=\small,anchor=north,yshift=-3pt] at (axis cs:2,-3.7) {$2$};

\end{axis}
\end{tikzpicture}
\par\small $f(\pi/3)=\sqrt3=2(1/\sqrt3)/(1-1/3)$。虚线是无定义位置，不是曲线的一部分。
\end{center}


由 $\tan(x+y)=\dfrac{\sin(x+y)}{\cos(x+y)}$，把两角和的正弦、余弦公式代入，并同除以 $\cos x\cos y$，得到
\[
\boxed{\tan(x+y)=\frac{\tan x+\tan y}{1-\tan x\tan y}}.
\]
\textbf{快记：}上面\textbf{相加}，下面\textbf{1 减积}。如抽象函数满足这一式，只要代入 $x=y$ 就可求 $f(2x)$。

\textbf{条件：}三个正切值必须有定义，而且分母不为 $0$；所以绝不能把公式当作任意实数都能代入的关系。
\end{explanationbox}

\FReserve{7.5cm}
\begin{explanationbox}
\textbf{\textcolor{CExplanation}{7 正弦型：和差相乘，平方相减}}
% F085｜真实数学图；代表函数不构成抽象解的唯一性证明。
\begin{center}
\begin{tikzpicture}
\begin{axis}[width=12.4cm,height=5.1cm,axis lines=middle,
 xmin=-0.3,xmax=6.8,ymin=-1.4,ymax=1.5,
 xtick={0,1.570796,3.141593,4.712389,6.283185},xticklabels=\empty,ytick={-1,0,1},xlabel={$t$},ylabel={$f(t)$},
 title={正弦曲线 $f(t)=\sin t$},title style={font=\small,yshift=3pt},
 tick label style={font=\small},label style={font=\small},
 ylabel style={at={(axis description cs:0,1)},anchor=south east,xshift=-18pt},
 clip=false,grid=major,grid style={black!8},enlargelimits=false,
 every axis plot/.append style={thick},samples=161]
\addplot[CExplanation,very thick,domain=0:6.28] {sin(deg(x))};
\draw[dashed,black!35] (axis cs:0.5235987755982988,0) |- (axis cs:0,0.5);
\draw[dashed,black!35] (axis cs:1.0471975511965976,0) |- (axis cs:0,0.8660254037844386);
\draw[dashed,black!35] (axis cs:1.5707963267948966,0) |- (axis cs:0,1);
\addplot[only marks,mark=*,mark size=2pt,orange!80!black] coordinates {(0.523598776,0.5) (1.04719755,0.866025404) (1.57079633,1)};

\node[font=\small,anchor=north,yshift=-3pt] at (axis cs:1.570796,-1.4) {$\pi/2$};
\node[font=\small,anchor=north,yshift=-3pt] at (axis cs:3.141593,-1.4) {$\pi$};
\node[font=\small,anchor=north,yshift=-3pt] at (axis cs:4.712389,-1.4) {$3\pi/2$};
\node[font=\small,anchor=north,yshift=-3pt] at (axis cs:6.283185,-1.4) {$2\pi$};

\end{axis}
\end{tikzpicture}
\par\small 在 $x+y=\pi/2$ 与 $x-y=\pi/6$ 读出两高度，乘积等于 $\sin^2x-\sin^2y$。
\end{center}


来自熟悉的三角恒等式
\[
\boxed{\sin(x+y)\sin(x-y)=\sin^2x-\sin^2y}.
\]
还可以从两角和差正弦公式相乘得到：
\[
(\sin x\cos y+\cos x\sin y)
(\sin x\cos y-\cos x\sin y)
=\sin^2x-\sin^2y.
\]
\textbf{最快的观察：}遇到 $f(m)f(n)$，先试
\[x=\frac{m+n}{2},\quad y=\frac{m-n}{2}.\]
\textbf{不能只认正弦：}因为 $(x+y)(x-y)=x^2-y^2$，函数 $f(x)=x$ 也满足这个模型。
\end{explanationbox}

\FReserve{7.5cm}
\begin{explanationbox}
\textbf{\textcolor{CExplanation}{8 余切型：上面乘积减 1}}
% F085｜真实数学图；代表函数不构成抽象解的唯一性证明。
\begin{center}
\begin{tikzpicture}
\begin{axis}[width=12.4cm,height=5.1cm,axis lines=middle,
 xmin=-1.5,xmax=3.1,ymin=-3.9,ymax=4.1,
 xtick={-1,0,1,2,3},xticklabels=\empty,ytick={-3,-1,0,1,3},xlabel={$t$},ylabel={$f(t)$},
 title={余切曲线 $f(t)=\cot t$},title style={font=\small,yshift=3pt},
 tick label style={font=\small},label style={font=\small},
 ylabel style={at={(axis description cs:0,1)},anchor=south east,xshift=-18pt},
 clip=false,grid=major,grid style={black!8},enlargelimits=false,
 every axis plot/.append style={thick},samples=161]
\addplot[CExplanation,very thick,domain=.27:2.85] {cos(deg(x))/sin(deg(x))};
\addplot[CExplanation,very thick,domain=-1.3:-.27] {cos(deg(x))/sin(deg(x))};
\draw[dashed,black!35] (axis cs:0.5235987755982988,0) |- (axis cs:0,1.7320508075688772);
\draw[dashed,black!35] (axis cs:1.0471975511965976,0) |- (axis cs:0,0.5773502691896258);
\addplot[only marks,mark=*,mark size=2pt,orange!80!black] coordinates {(0.523598776,1.73205081) (1.04719755,0.577350269)};

\node[font=\small,anchor=north,yshift=-3pt] at (axis cs:-1,-3.9) {$-1$};
\node[font=\small,anchor=north,yshift=-3pt] at (axis cs:1,-3.9) {$1$};
\node[font=\small,anchor=north,yshift=-3pt] at (axis cs:2,-3.9) {$2$};
\node[font=\small,anchor=north,yshift=-3pt] at (axis cs:3,-3.9) {$3$};

\end{axis}
\end{tikzpicture}
\par\small 两次输入 $\pi/6$ 得 $\pi/3$，高度从 $\sqrt3$ 变为 $(3-1)/(2\sqrt3)=1/\sqrt3$；$t=0$ 无定义。
\end{center}


把 $\cot(x+y)=\dfrac{\cos(x+y)}{\sin(x+y)}$ 中的两角和公式代入，再同除以 $\sin x\sin y$，得到
\[
\boxed{\cot(x+y)=\frac{\cot x\cot y-1}{\cot x+\cot y}}.
\]
\textbf{快记：}上面\textbf{乘积减 1}，下面\textbf{相加}。它与正切模型的分子、分母恰好相互调换了重点。

\textbf{特别注意：}$\cot0$ 没有定义！因此不能像指数或正比例型一样直接令输入为 $0$。
\end{explanationbox}

